\documentclass[11pt]{article}
\usepackage[a4paper,margin=26mm]{geometry}
\usepackage{amsmath,amssymb,mathtools}
\usepackage{unicode-math}
\setmainfont{Latin Modern Roman}
\setmathfont{Latin Modern Math}
\usepackage{xurl,hyperref}
\hypersetup{hidelinks,pdftitle={Problem statement}}
\newcommand{\sref}[2]{\hyperlink{source:#1:#2}{[S#2]}}
\newcommand{\cataloguescope}[1]{#1}
\begin{document}
\section{Coarse foliations — Question 12.2}\label{statement}

\subsection{Definitions and mathematical statement}

\paragraph{Shared notation.}$\mathbb N=\{1,2,\ldots\}$, and $\mathbb Z,\mathbb Q,\mathbb R,\mathbb C$ denote the integers, rationals, reals and complex numbers. Any other conventions are specified locally.


A cone field on a smooth $3$-manifold $M$ is a smooth symmetric covariant $2$-tensor $C$ of signature $(2,1)$. An orientation of $C$ is a continuous choice at every $p\in M$ of one of the two components of $\{v\in T_pM:C_p(v,v)<0\}$, called the positive timelike cone. A positively timelike path has its tangent in that cone. The cone field is recurrent if every ordered pair of points of $M$ is joined by such a path. It is homotopically essential if every closed positively timelike path is noncontractible in $M$. A taut cone field is an oriented cone field with both properties. These are the conventions of Definition 2.4 in [S3]. A supported embedded loop is an embedded closed positively timelike path. For such a loop $K$, Dehn surgery removes a tubular neighborhood of $K$ and attaches a solid torus whose meridian is identified with a slope, an unoriented primitive curve class on the boundary torus. Nontrivial surgery means a slope different from the original meridian.
\paragraph{Statement recorded in the supplied source.}
Does every hyperbolic $3$-manifold admit a taut cone field? For a taut cone field and a supported embedded loop, can the cone field remain taut after surgery at every slope outside a finite set? Is there a class of examples in which it remains taut after every nontrivial surgery? The surgery questions include constructing a suitable extension of the cone field over the filled solid torus while retaining it away from the surgery neighborhood. \sref{statement}{2}
\subsection{Short English statement}
Does every hyperbolic three-manifold admit a recurrent cone field with only noncontractible supported loops, and how broadly can tautness survive surgery on a supported loop?
\subsection{Sources}
\begin{description}
\item[\hypertarget{source:statement:1}{S1}] ULAM.AI and the UnsolvedMath Contributors, \emph{UnsolvedMath: A Curated Collection of Open Mathematics Problems}, catalogue entry. CC BY 4.0. \url{https://huggingface.co/datasets/ulamai/UnsolvedMath}.
\item[\hypertarget{source:statement:2}{S2}] Danny Calegari, \emph{Problems in foliations and laminations of 3-manifolds}, version 0.78 (2002), Question 12.2; Definitions 1.1--1.10 and local definitions and Remark (2). \url{https://arxiv.org/abs/math/0209081v1}.
\item[\hypertarget{source:statement:3}{S3}] Danny Calegari, \emph{Useful branched surfaces which carry nothing}, version 0.82, arXiv:math/0010146v1 (2000), Definition 2.4, pp. 5--6 (withdrawn preprint, cited here for its cone-field terminology). \url{https://arxiv.org/abs/math/0010146v1}.
\end{description}
\label{end:statement}
\end{document}
